\documentclass[12pt]{article}
\usepackage[margin=1in]{geometry}
\title{Long-Form Prose}
\author{Technical Writer}
\begin{document}
\maketitle

\section{The Challenge of PDF Conversion}
Converting PDF documents to Markdown is fundamentally challenging because the two
formats represent content in fundamentally different ways. PDF is a page description
language that specifies the exact position, font, and size of each character on a
page. It is designed for precise visual reproduction across different devices and
platforms, prioritizing appearance over structure. Markdown, on the other hand, is
a semantic markup language that describes the logical structure of content through
simple syntax conventions like hash marks for headings, asterisks for emphasis,
and pipe characters for tables.

The gap between these two paradigms means that any converter must bridge the
divide between visual presentation and logical structure. This requires the
converter to reverse-engineer the author's intent from the visual layout, inferring
paragraph boundaries from spacing, heading levels from font sizes, and list items
from indentation and marker characters. This process is inherently heuristic and
imperfect, because the same visual appearance can result from different logical
structures, and different tools may produce subtly different PDF output for the
same source document.

\section{Our Approach}
Our six-layer pipeline approaches this challenge by separating concerns into
distinct processing stages. The first layer extracts raw character data from the
PDF, including precise position coordinates, font information, and text content.
The second layer assembles these characters into words, lines, and blocks using
adaptive gap analysis that adjusts thresholds based on the local font size. The
third layer performs layout analysis to detect headers, footers, columns, and
reading order. The fourth layer applies semantic classification to identify
headings, paragraphs, lists, and other structural elements. The fifth layer
handles table extraction using specialized algorithms. The sixth and final layer
generates the Markdown output.

\section{Content Accuracy as the Primary Goal}
Throughout the entire pipeline, our primary goal is content accuracy rather than
formatting fidelity. We prioritize ensuring that every word, number, and symbol
from the source PDF appears in the output Markdown, even if the formatting is
not perfectly replicated. This design decision reflects the intended use case:
converting documents for consumption by AI and language models, where content
completeness matters more than visual appearance.
\end{document}